\section{Reproducing this report}
\label{app:repro}

\begin{table}[htbp]
\centering
\caption{Everything in this document is generated. The order below is the order
to run it in. The fitting step dominates; the rest takes a few minutes.}
\label{tab:repro}
\small
\begin{tabular}{@{}p{5.4cm}p{9.4cm}@{}}
\toprule
\textbf{Command} & \textbf{What it produces} \\
\midrule
\code{pip install -r requirements-dev.txt} & The full analysis environment.
\code{requirements.txt} alone is the deployment runtime \\[3pt]
\code{python tests/test\_leakage.py} & \checkcount{} validation checks: join
grain, the history construction against a brute-force recomputation, the leakage
guard in both directions, split disjointness, Phase 1 backward compatibility,
the capacity and lateness tables against recomputation, and --- added in this
phase --- that the pickled models, the metadata, the hyperparameter table, the
README and this report name the same deployed estimators, and that the
review-timing value recomputes from the per-order predictions \\[3pt]
\code{python src/train\_models.py} & Fits every model, scores the test window
once, writes \code{models/} and every result table into \code{report/results/} \\[3pt]
\code{python src/capacity\_metrics.py} & Precision@$K$ and Recall@$K$ from the
saved per-order test predictions (Section~\ref{sec:results-capacity}); nothing
is refitted \\[3pt]
\code{python src/lateness\_anatomy.py} & The weekly, monthly and strike-cohort
lateness tables of Appendix~\ref{app:lateness} \\[3pt]
\code{python src/fit\_diagnostics.py} & The three deployed models scored on the
training, validation and test periods (Appendix~\ref{app:tvt}) \\[3pt]
\code{python src/selection\_tables.py} & Every candidate's score on every
out-of-sample window with its fold-to-fold spread (Appendix~\ref{app:protocol}),
and the selected hyperparameters of every searched model and of the deployed
ensembles' members, read from the pickles (Appendix~\ref{app:hyperparameters}) \\[3pt]
\code{python src/pilot\_design.py} & The sizing of the outreach pilot
(Section~\ref{sec:pilot}, Appendix~\ref{app:pilot}) \\[3pt]
\code{python src/review\_timing.py} & When each review was written against the
delivery, the share of the call list that can still be reached, the value
restated on that share, and the per-order flags the application's replay shows
(Appendix~\ref{app:timing}) \\[3pt]
\code{python src/figures\_phase2.py} & Draws every modelling figure as vector
PDF into \code{report/figures/} from those tables; the exploratory figures are
drawn by notebooks 01--04 \\[3pt]
\code{python src/phase3\_stats.py} & Distils every number quoted in this report
into \code{report/stats\_phase3.txt} (it runs \code{phase2\_stats.py} first, so
the Phase 2 contract is a subset). Diff it against the manuscript \\[3pt]
\code{python src/report\_stats.py} & Reproduces every exploratory number of
Section~\ref{sec:eda} into \code{report/stats.txt}, unchanged since Phase 1 ---
the check that the later feature additions are backward compatible \\[3pt]
\code{streamlit run deployment/app.py} & The application of
Section~\ref{sec:deployment}; \code{dashboard/app.py} is the Community Cloud
entry point and runs the same file \\
\bottomrule
\end{tabular}
\end{table}

Table~\ref{tab:repro} is the run order. \code{SEED = 42} is set in every
estimator, every randomised search and every subsample. Both requirements files
pin exact versions on Python $\le$ 3.12, the
interpreter every reported number was produced on. Model artifacts record the
\code{scikit-learn} version that wrote them and \code{modeling.load\_model()}
checks it on the way back in. The nine notebooks (01--04 exploratory, 05--09
modelling, 10 the additions of this phase) are committed with their outputs.

\section{Data model, cleaning decisions and dictionary}
\label{app:dictionary}

\begin{figure}[htbp]
\centering
\resizebox{\linewidth}{!}{%
\begin{tikzpicture}[
  font=\small,
  box/.style={draw, rounded corners=2pt, minimum height=9mm, minimum width=25mm,
              inner xsep=4pt, inner ysep=3pt, align=center},
  spine/.style={box, very thick},
  link/.style={draw, -{Latex[length=4pt]}, thin},
  lbl/.style={font=\scriptsize, midway, fill=white, inner sep=1pt}
]
\node[spine] (ord) at (0,0) {\textbf{ORDERS}\\[-2pt]\scriptsize 99{,}441\\[-2pt]\scriptsize \textit{one row = one order}};

\node[box] (cus) at (-4.3,1.5) {\textbf{CUSTOMERS}\\[-2pt]\scriptsize 99{,}441};
\node[box] (itm) at (4.3,1.5) {\textbf{ORDER\_ITEMS}\\[-2pt]\scriptsize 112{,}650};
\node[box] (pay) at (4.3,0) {\textbf{ORDER\_PAYMENTS}\\[-2pt]\scriptsize 103{,}886};
\node[box] (rev) at (4.3,-1.5) {\textbf{ORDER\_REVIEWS}\\[-2pt]\scriptsize 99{,}224};

\node[box] (prd) at (8.6,2.2) {\textbf{PRODUCTS}\\[-2pt]\scriptsize 32{,}951};
\node[box] (sel) at (8.6,0.8) {\textbf{SELLERS}\\[-2pt]\scriptsize 3{,}095};
\node[box] (cat) at (12.4,2.2) {\textbf{CATEGORY}\\[-2pt]\scriptsize \textbf{TRANSLATION} 71};
\node[box] (geo) at (-4.3,-1.0) {\textbf{GEOLOCATION}\\[-2pt]\scriptsize 1{,}000{,}163};

\draw[link] (cus) -- node[lbl]{1\,:\,1} (ord);
\draw[link] (ord) -- node[lbl]{1\,:\,N} (itm);
\draw[link] (ord) -- node[lbl]{1\,:\,N} (pay);
\draw[link] (ord) -- node[lbl]{1\,:\,N} (rev);
\draw[link] (itm) -- node[lbl]{N\,:\,1} (prd);
\draw[link] (itm) -- node[lbl]{N\,:\,1} (sel);
\draw[link] (prd) -- node[lbl]{N\,:\,1} (cat);
\draw[link] (geo) -- node[lbl,align=center]{zip\\prefix} (cus);
\end{tikzpicture}}
\caption{The nine source tables. \code{ORDERS} is the analysis grain. The three
one-to-many children (items, payments and reviews) are each aggregated to one
row per order \emph{before} joining; merging them directly inflates the table by
19\% and biases every average.}
\label{fig:datamodel}
\end{figure}

\begin{table}[htbp]
\centering
\caption{Data quality audit and cleaning decisions, implemented as
\code{CLEANING\_DECISIONS} in \code{src/data\_load.py}.}
\label{tab:quality}
\small
\begin{tabular}{@{}p{2.7cm}p{4.7cm}p{6.3cm}@{}}
\toprule
\textbf{Issue} & \textbf{Evidence} & \textbf{Decision} \\
\midrule
One-to-many joins &
112{,}650 items and 103{,}886 payments against 99{,}441 orders; a naive merge
yields 118{,}437 rows &
Aggregate each child table to one row per order, then join; assert the grain \\[3pt]
Customer key &
99{,}441 \code{customer\_id} values but 96{,}096 \code{customer\_unique\_id};
2{,}997 people ordered more than once &
Carry \code{customer\_unique\_id} and group on it for repeat behaviour and for
splitting \\[3pt]
Undelivered orders &
2.98\% of orders have no delivery timestamp; status is \emph{delivered} for
97.0\% &
Delivery targets computed only for delivered orders; the rest are \code{NaN},
never imputed \\[3pt]
Review keys &
\code{review\_id} is not unique (789 ids reused); 547 orders carry 2--3 reviews
and 768 carry none &
Key on \code{order\_id}, keep the earliest review; no review means \code{NaN} \\[3pt]
Sparse review text &
\code{review\_comment\_title} present for 11.7\% of reviews,
\code{review\_comment\_message} for 41.3\% &
Use \code{review\_score} only; coverage does not support NLP features \\[3pt]
Geolocation grain &
1{,}000{,}163 rows for 19{,}010 zip prefixes, including points outside Brazil &
Collapse to a median coordinate per prefix; use it for great-circle distance \\[3pt]
Missing product attributes &
610 products (1.9\%) have no category &
Label \code{unknown} rather than dropping orders; impute numerics in-pipeline \\[3pt]
Coverage ramp &
326 orders before 2017; final weeks truncated mid-cycle &
Restrict to 2017-01-01 .. 2018-08-31 (99{,}091 orders) \\[3pt]
Portuguese categories &
71 category names &
Translate to English for all output; keep the original key \\
\bottomrule
\end{tabular}
\end{table}

\begin{table}[htbp]
\centering
\caption{Candidate problems against the brief's scoping checklist
(Section~\ref{sec:candidates}), as scored in Phase 1 with untuned probes.}
\label{tab:checklist}
\small
\begin{tabular}{@{}p{2.5cm}p{4.1cm}p{3.6cm}p{3.5cm}@{}}
\toprule
\textbf{Criterion} & \textbf{C1 delivery (checkout)} & \textbf{C2 low review (delivery)} & \textbf{C3 repeat (checkout)} \\
\midrule
Derivable target & Yes, from order timestamps & Yes, \code{review\_score} $\le 2$ & Yes, next order by person \\
No leakage & Yes, checkout regime & Yes, at-delivery regime & Yes, checkout regime \\
Sufficient signal & Regression yes ($-41\%$ MAE); classification weak (1.6$\times$) & Yes, PR-AUC 0.40 vs 0.11 (3.6$\times$) & No, 2.0$\times$ at a 1.7\% base rate \\
Manageable imbalance & 8.1\% late; PR-AUC and thresholds, no resampling & 14.6\% low; same handling & 2.6\%, too extreme \\
Clear stakeholder & Operations & Customer experience & CRM \\
\midrule
\textbf{Verdict} & \textbf{Problem 1} & \textbf{Problem 2} & Rejected \\
\bottomrule
\end{tabular}
\end{table}

\begin{table}[htbp]
\centering
\caption{Data dictionary for the nine source tables. \emph{Null} is measured,
not declared: Y means the column actually contains missing values in this
extract, with the observed rate given.}
\label{tab:dictionary}
\scriptsize
\setlength{\tabcolsep}{3.5pt}
\renewcommand{\code}[1]{\texttt{#1}}
\begin{tabular}{@{}p{2.0cm}p{4.4cm}p{1.1cm}p{1.05cm}p{5.5cm}@{}}
\toprule
\textbf{Table} & \textbf{Column} & \textbf{Type} & \textbf{Null} & \textbf{Comment} \\
\midrule
\code{orders} & \code{order\_id} & text & N & Primary key. The analysis grain \\
 & \code{customer\_id} & text & N & \textbf{Per-order} key, not a person \\
 & \code{order\_status} & text & N & 8 values; \emph{delivered} for 97.0\% \\
 & \code{order\_purchase\_timestamp} & datetime & N & The checkout prediction moment \\
 & \code{order\_approved\_at} & datetime & Y 0.16\% & Post-checkout; excluded as a feature \\
 & \code{order\_delivered\_carrier\_date} & datetime & Y 1.79\% & Post-outcome; splits lead time into seller and carrier share \\
 & \code{order\_delivered\_customer\_date} & datetime & Y 2.98\% & Post-outcome; defines both delivery targets \\
 & \code{order\_estimated\_delivery\_date} & date & N & Promised at checkout; \textbf{admissible}. Date only, 459 distinct values \\
\midrule
\code{order\_items} & \code{order\_id} & text & N & \textbf{1:N} --- 112,650 rows for 98,666 orders \\
 & \code{order\_item\_id} & int & N & Sequence within the order, 1--21 \\
 & \code{product\_id}, \code{seller\_id} & text & N & Foreign keys \\
 & \code{shipping\_limit\_date} & datetime & N & Post-checkout; excluded \\
 & \code{price}, \code{freight\_value} & float & N & Per item; summed to order level \\
\midrule
\code{order\_payments} & \code{order\_id} & text & N & \textbf{1:N}. One order has no payment row \\
 & \code{payment\_sequential} & int & N & 1--29 \\
 & \code{payment\_type} & text & N & 5 values; credit card 75.4\% \\
 & \code{payment\_installments} & int & N & 0--24; $>1$ for 51.4\% of orders \\
 & \code{payment\_value} & float & N & Summed to order level \\
\midrule
\code{order\_reviews} & \code{review\_id} & text & N & \textbf{Not unique} --- 789 ids reused across orders \\
 & \code{order\_id} & text & N & \textbf{1:N}; 547 orders carry 2--3 reviews \\
 & \code{review\_score} & int & N & 1--5; target of Problem 2 \\
 & \code{review\_comment\_title} & text & Y 88.3\% & Too sparse to use \\
 & \code{review\_comment\_message} & text & Y 58.7\% & Too sparse to use; no NLP features \\
 & \code{review\_creation\_date} & date & N & Post-outcome for the review target \\
 & \code{review\_answer\_timestamp} & datetime & N & Post-outcome \\
\midrule
\code{customers} & \code{customer\_id} & text & N & Joins 1:1 to orders \\
 & \code{customer\_unique\_id} & text & N & \textbf{The actual person}; 96,096 of them. Grouping key for splits \\
 & \code{customer\_zip\_code\_prefix} & int & N & Joins to geolocation \\
 & \code{customer\_city}, \code{customer\_state} & text & N & 4,119 cities, 27 states \\
\midrule
\code{sellers} & \code{seller\_id} & text & N & 3,095 sellers \\
 & \code{seller\_zip\_code\_prefix} & int & N & Joins to geolocation \\
 & \code{seller\_city}, \code{seller\_state} & text & N & 59.7\% of sellers are in SP \\
\midrule
\code{products} & \code{product\_id} & text & N & 32,951 products \\
 & \code{product\_category\_name} & text & Y 1.85\% & Portuguese; 73 values. Nulls labelled \emph{unknown} \\
 & \code{product\_weight\_g} & float & Y 0.01\% & Summed to order level \\
 & \code{product\_length/height/width\_cm} & float & Y 0.01\% & Combined into item volume \\
 & \code{product\_photos\_qty} & float & Y 1.85\% & Listing-quality proxy \\
\midrule
\code{geolocation} & \code{geolocation\_zip\_code\_prefix} & int & N & 19,015 prefixes over 1,000,163 rows \\
 & \code{geolocation\_lat}, \code{\_lng} & float & N & Collapsed to a median per prefix; out-of-Brazil points dropped \\
\midrule
\code{category\_} & \code{product\_category\_name} & text & N & Join key \\
\code{translation} & \code{product\_category\_name\_english} & text & N & 71 translations \\
\midrule
\multicolumn{5}{@{}l@{}}{\textit{Derived in \code{src/} --- not present in any source file:}} \\
derived & \code{delivery\_days} & float & Y & Purchase to arrival; target of 1A \\
derived & \code{is\_late} & bool & Y & Arrival after the promised date; target of 1B \\
derived & \code{distance\_km} & float & Y 1.2\% & Great-circle customer--seller distance \\
derived & \code{estimated\_days} & float & N & Promised window in days \\
derived & \code{lead\_category} & text & N & Category of the most expensive item \\
derived & \code{freight\_ratio} & float & Y & Freight as a share of item value \\
derived & \code{hist\_*} & float & Y & The eight as-of track records of Section~\ref{sec:history} \\
\bottomrule
\end{tabular}
\end{table}

\clearpage
\section{Success criteria registered in advance}
\label{app:bands}

Written into \code{config.EXPECTED\_BANDS} before any model was tuned.
\code{evaluate.check\_bands()} compares the headline results against them at the
end of every training run and writes its verdicts to
\code{report/results/bands.csv}. The alarm is direction-aware: for an error
metric the suspicious end of the interval is the low one, for a score the high
one, and confusing the two would point the check at the shortfalls and wave the
leaks through.

\begin{center}
\small
\begin{tabular}{@{}llp{2.4cm}p{5.0cm}r@{}}
\toprule
\textbf{Problem} & \textbf{Metric} & \textbf{Registered band} & \textbf{Where the range came from} & \textbf{Result} \\
\midrule
1A & MAE (days) & 3.3 -- 4.3 & Phase 1 reported 3.89 untuned & 3.898 \\
1B & balanced accuracy & 0.58 -- 0.72 & Phase 1 reported 0.58 for a direct classifier & 0.705 \\
1B & ROC-AUC & 0.72 -- 0.85 & Untuned scoping probe reached 0.79 for the derived route & 0.782 \\
2 & balanced accuracy & 0.63 -- 0.74 & Phase 1 reported 0.69 & 0.681 \\
2 & ROC-AUC & 0.70 -- 0.79 & Phase 1 reported 0.73 & 0.749 \\
2 & PR-AUC & 0.34 -- 0.48 & Phase 1 reported 0.40 against a 0.11 floor & 0.414 \\
\bottomrule
\end{tabular}
\end{center}

\section{Hyperparameter search spaces and selected values}
\label{app:hyperparameters}

\begin{center}
\scriptsize
\begin{tabular}{@{}llll@{}}
\toprule
\textbf{Family} & \textbf{Rung} & \textbf{Classification} & \textbf{Regression} \\
\midrule
A --- Linear     & 0 & \code{LogisticRegression} & \code{LinearRegression} \\
                 & 1 & \code{LogisticRegression} ($\ell_2$, tuned $C$) & \code{Ridge} (tuned $\alpha$) \\
B --- Tree-based & 0 & \code{DecisionTreeClassifier} & \code{DecisionTreeRegressor} \\
                 & 1 & \code{RandomForestClassifier} & \code{RandomForestRegressor} \\
                 & 2 & \code{HistGradientBoostingClassifier} & \code{HistGradientBoostingRegressor} \\
C --- Ensemble   & 0 & \code{VotingClassifier} (soft, A+B) & \code{VotingRegressor} (A+B) \\
                 & 1 & \code{StackingClassifier} (A+B) & \code{StackingRegressor} (A+B) \\
\bottomrule
\end{tabular}
\end{center}

Every tuned model was searched with \code{RandomizedSearchCV} over the
distributions of Table~\ref{tab:spaces}, scored by five-fold expanding-window
cross-validation inside the training period: thirty draws per model, except the
random forests, which get six because a forest costs one to two orders of
magnitude more per fit and Phase 1 had already shown the forest and the boosted
trees landing within 0.01 days of each other. Table~\ref{tab:selected} lists the
values that won, which the Phase 2 report referred to but did not print; the
deployed ensembles carry exactly their winning members' values, read back from
the pickles, and the validation suite asserts it.

\begin{table}[htbp]
\centering
\caption{Search spaces.}
\label{tab:spaces}
\small
\begin{tabular}{@{}p{4.1cm}cp{7.9cm}@{}}
\toprule
\textbf{Model} & \textbf{Draws} & \textbf{Distributions} \\
\midrule
\code{LinearRegression} / \code{LogisticRegression} & --- &
No free parameters. Fitted and scored under the same protocol so the row is
comparable \\[3pt]
\code{Ridge} & 30 & $\alpha \sim \mathrm{loguniform}(10^{-3}, 10^{3})$ \\[3pt]
\code{LogisticRegression} ($\ell_2$) & 30 & $C \sim \mathrm{loguniform}(10^{-3}, 10^{2})$ \\[3pt]
\code{DecisionTree} & 30 & \code{max\_depth} $\sim$ randint(3, 20) [classifier: (3, 16)];
\code{min\_samples\_leaf} $\sim$ randint(5, 200) [classifier: (20, 400)];
\code{max\_features} $\sim U(0.4, 1.0)$ \\[3pt]
\code{RandomForest} & 6 & 150 trees; \code{max\_depth} $\sim$ randint(8, 18)
[classifier: (6, 16)]; \code{min\_samples\_leaf} $\sim$ randint(20, 140)
[classifier: (30, 180)]; \code{max\_features} $\sim U(0.15, 0.40)$ \\[3pt]
\code{HistGradientBoosting} & 30 & \code{max\_iter} $\sim$ randint(120, 500);
\code{learning\_rate} $\sim \mathrm{loguniform}(0.02, 0.3)$;
\code{max\_leaf\_nodes} $\sim$ randint(15, 64);
\code{min\_samples\_leaf} $\sim$ randint(10, 120);
\code{l2\_regularization} $\sim \mathrm{loguniform}(10^{-4}, 10)$ \\[3pt]
\code{Voting} & --- & Weights fixed at 1:2 in the tree's favour. A stated prior,
not a tuned parameter: under drift the linear member is materially weaker and an
equal blend would be worse than its better half \\[3pt]
\code{Stacking} & --- & Meta-learner \code{Ridge}($\alpha=1$) /
\code{LogisticRegression}($C=1$); inner CV five contiguous folds --- see the note below \\
\bottomrule
\end{tabular}
\end{table}

\begin{table}[htbp]
\centering
\caption{Selected hyperparameter values, from
\code{report/results/hyperparameters.csv}. Members of the deployed ensembles are
the rows marked $\dagger$; the stacked review model's meta-learner has
coefficients 0.405 (linear member) and 5.572 (tree member).}
\label{tab:selected}
\footnotesize
\setlength{\tabcolsep}{4pt}
\begin{tabular}{@{}llp{9.6cm}@{}}
\toprule
\textbf{Problem} & \textbf{Model} & \textbf{Selected values} \\
\midrule
1A & LinearRegression $\dagger$ & no free parameters \\
1A & Ridge & $\alpha = 659.9$ \\
1A & DecisionTree & depth 17, min leaf 158, max features 0.897 \\
1A & RandomForest & 150 trees, depth 14, min leaf 107, max features 0.165 \\
1A & HistGradientBoosting $\dagger$ & 433 iterations, learning rate 0.1065, 36 leaf nodes, min leaf 117, $\ell_2$ 9.142 \\
1A & VotingRegressor $\dagger$ & weights 1:2 (linear : tree) \\
\midrule
1B & LogisticRegression ($\ell_2$) & $C = 0.004983$ \\
1B & DecisionTree & depth 9, min leaf 290, max features 0.878 \\
1B & RandomForest & 150 trees, depth 13, min leaf 132, max features 0.300 \\
1B & HistGradientBoosting & 426 iterations, learning rate 0.07079, 21 leaf nodes, min leaf 30, $\ell_2$ 8.244 \\
1B & derived route $\dagger$ & one-parameter logistic on the margin, fitted on validation; threshold 0.11 \\
\midrule
2 & LogisticRegression ($\ell_2$) $\dagger$ & $C = 8.431$ \\
2 & DecisionTree & depth 14, min leaf 211, max features 0.963 \\
2 & RandomForest & 150 trees, depth 12, min leaf 44, max features 0.349 \\
2 & HistGradientBoosting $\dagger$ & 208 iterations, learning rate 0.02073, 28 leaf nodes, min leaf 18, $\ell_2$ 0.01789 \\
2 & StackingClassifier $\dagger$ & meta-learner LogisticRegression($C = 1$); threshold 0.19 \\
\bottomrule
\end{tabular}
\end{table}

\paragraph{A compromise inside stacking, stated plainly.}
\code{StackingClassifier} and \code{StackingRegressor} build their meta-features
with \code{cross\_val\_\allowbreak predict}, which requires the inner folds to \emph{partition}
the training set: every row must be predicted exactly once. An expanding window
partitions nothing --- the earliest rows are never in a test fold --- so the
closest valid approximation is unshuffled \code{KFold}, which at least keeps each
fold contiguous in time. The meta-learner is consequently fitted on out-of-fold
predictions partly made by models trained on later orders. The residual optimism
is confined to the meta-learner's coefficients; the base models are refitted on
the whole training period afterwards, and the outer protocol scoring the
finished ensemble is still the expanding window. It matters less than it
otherwise would because of how the result lands: the ensemble is
\emph{flattered} by this arrangement and still does not separate from the
boosted trees, which makes the conclusion in
Section~\ref{sec:results-ensemble} safer rather than weaker.

\section{Re-checking a Phase 1 claim about least squares}
\label{app:encoding}

Phase 1 reported that ordinary least squares was unusable on this design matrix
--- ``rank-deficient, so its coefficients explode''. Phase 2 re-ran it four ways.
The observation reproduces exactly; the conclusion drawn from it needed
correcting.

\begin{table}[htbp]
\centering
\caption{Reproducible via \code{report/results/encoding\_diagnostic.csv}. The
deficiency is exactly four --- one per categorical, because each one-hot block
sums to one and is therefore collinear with the intercept.}
\label{tab:encoding}
\small
\begin{tabular}{@{}lccrrr@{}}
\toprule
\textbf{Encoding} & \textbf{rank / cols} & \textbf{deficiency} &
\textbf{max $|\beta|$} & \textbf{CV MAE} & \textbf{test MAE} \\
\midrule
OLS, full one-hot block   & 101 / 105 & 4 & $2.7\times10^{13}$ & $3.9\times10^{11}$ & 4.181 \\
OLS, one level dropped    & 101 / 101 & 0 & 14.7  & 5.242 & 4.181 \\
Ridge, full one-hot block & 101 / 105 & 4 & 9.82  & 5.241 & 4.182 \\
Ridge, one level dropped  & 101 / 101 & 0 & 12.5  & 5.240 & 4.183 \\
\bottomrule
\end{tabular}
\end{table}

Least squares still returns a minimum-norm solution and its \emph{predictions}
are fine --- 4.181 days on the test window, indistinguishable from Ridge. What is
not fine is the coefficient vector at $10^{13}$, and the cancellation holding it
together is fragile enough to fail on one of five expanding-window folds, which
is why the cross-validated error is astronomical while the single-split test
error looks healthy. Dropping one level per categorical restores full rank, so
the linear family is encoded that way; the compromise is that a level unseen in
training is scored as the reference level, which affects a handful of test
orders with a payment type absent from training. Once the design is full rank,
Ridge and OLS agree to three decimals: regularisation buys numerical stability
here, not accuracy.

\section{Full result tables}
\label{app:tables}

The body reports the models that matter. The complete tables --- every model,
every metric, cross-validated, validation and test scores side by side --- are
written to \code{report/results/}:

\begin{center}
\footnotesize
\begin{tabular}{@{}p{5.2cm}p{9.6cm}@{}}
\toprule
\textbf{File} & \textbf{Contents} \\
\midrule
\code{reg\_p1a.csv} & Problem 1A: baselines, both ladders, ensembles; MAE, RMSE, $R^2$ \\
\code{clf\_p1b.csv} & Problem 1B: five direct classifiers and two derived routes \\
\code{clf\_p2.csv} & Problem 2: both ladders and both ensembles, all metrics \\
\code{selection\_*.csv}, \code{cv\_folds.csv} & Every candidate on every out-of-sample window, its ranks, and the fold-to-fold spread (Appendix~\ref{app:protocol}) \\
\code{hyperparameters.csv} & Selected values of every searched model and of the deployed ensembles' members (Appendix~\ref{app:hyperparameters}) \\
\code{encoding\_diagnostic.csv} & The Phase 1 re-check of Table~\ref{tab:encoding} \\
\code{cv\_protocol.csv} & Temporal against shuffled CV against the forward test \\
\code{weight\_ablation.csv} & Class weighting: discrimination and calibration \\
\code{calibration\_p2.csv}, \code{threshold\_sweep\_p2.csv} & Reliability and the operating-point sweep \\
\code{cost\_p2.csv}, \code{cost\_p1b.csv} & Expected net benefit by threshold \\
\code{perm\_p1a.csv}, \code{perm\_p2.csv} & Permutation importance on the test window \\
\code{segments\_*.csv} & Performance by region, category and test month \\
\code{predictions\_*.csv} & Per-order test predictions, from which the figures and the application's replay are drawn \\
\code{bands.csv} & Headline metrics against the bands registered in advance \\
\code{capacity\_*.csv} & Precision@$K$ and Recall@$K$, pooled and per day, for both flags \\
\code{lateness\_*.csv} & The lateness anatomy of Appendix~\ref{app:lateness} \\
\code{train\_val\_test.csv} & The deployed models on the training, validation and test periods (Appendix~\ref{app:tvt}) \\
\code{pilot\_design.csv}, \code{pilot\_breakeven.csv} & The outreach pilot sizing (Appendix~\ref{app:pilot}) \\
\code{review\_timing*.csv}, \code{story\_order.csv} & Review timing against delivery, the reachable share of the call list, the restated value, the on-time check and the example order (Appendix~\ref{app:timing}) \\
\bottomrule
\end{tabular}
\end{center}

\section{The cross-validation protocol, and every candidate on every window}
\label{app:protocol}

\begin{table}[htbp]
\centering
\caption{One model, three protocols. Raw average precision disagrees; the same
numbers divided by each period's no-skill floor do not.
Figure~\ref{fig:protocol}.}
\label{tab:protocol}
\small
\begin{tabular}{@{}lrrrr@{}}
\toprule
\textbf{Protocol} & \textbf{PR-AUC} & \textbf{fold spread} & \textbf{no-skill floor} & \textbf{lift} \\
\midrule
\code{TimeSeriesSplit}, expanding window & 0.562 & $\pm$0.091 & 0.158 & 3.55$\times$ \\
\code{StratifiedGroupKFold}, shuffled    & 0.580 & $\pm$0.006 & 0.158 & 3.66$\times$ \\
Forward test window --- the truth        & 0.414 & ---        & 0.110 & 3.75$\times$ \\
\bottomrule
\end{tabular}
\end{table}

Both cross-validated estimates in Table~\ref{tab:protocol} sit about 0.16 above
the forward result, which looks like catastrophic optimism --- and almost all
of it is mechanical. Average
precision's no-skill baseline \emph{is} the positive rate
\citep{saito2015prplot}, and the low-review rate falls from 15.8\% inside the
training period to 11.0\% in the test window; normalised by that floor the three
agree to within 0.2 of a multiple. \textbf{Raw PR-AUC is not comparable across
periods with different base rates; its lift over the floor is} --- and quoting
the raw number alone would have supported either a claim of severe overfitting or
a claim of none. The shuffled protocol adds a second warning: its fold-to-fold
spread is $\pm$0.006 against the temporal protocol's $\pm$0.091, an interval
fifteen times narrower on an estimate no closer to the truth.

\begin{table}[htbp]
\centering
\caption{Every Problem 2 candidate on every out-of-sample window (PR-AUC), from
\code{report/results/cv\_folds.csv}; the standard deviation is over the five
expanding folds. ``Rank'' is the mean rank over the three most recent windows
(folds 4, 5 and validation), the selection criterion of
Section~\ref{sec:cv}. Figure~\ref{fig:cvfolds}. The same table for Problems 1A
and 1B is in the CSV.}
\label{tab:cvfolds}
\scriptsize
\setlength{\tabcolsep}{2.6pt}
\resizebox{\linewidth}{!}{%
\begin{tabular}{@{}lrrrrrrrrrr@{}}
\toprule
\textbf{Model} & \textbf{fold 1} & \textbf{fold 2} & \textbf{fold 3} & \textbf{fold 4} & \textbf{fold 5} & \textbf{val.} & \textbf{mean $\pm$ sd} & \textbf{test} & \textbf{rank} \\
\midrule
StackingClassifier (A+B) \emph{deployed} & 0.413 & 0.537 & 0.617 & 0.660 & 0.614 & 0.409 & 0.568 $\pm$ 0.087 & 0.414 & 1.67 \\
HistGradientBoosting   & 0.411 & 0.539 & 0.618 & 0.660 & 0.613 & 0.407 & 0.568 $\pm$ 0.088 & 0.415 & 2.00 \\
RandomForest           & 0.404 & 0.544 & 0.622 & 0.654 & 0.616 & 0.403 & 0.568 $\pm$ 0.089 & 0.405 & 2.67 \\
VotingClassifier (A+B) & 0.413 & 0.526 & 0.611 & 0.650 & 0.604 & 0.404 & 0.561 $\pm$ 0.084 & 0.407 & 3.67 \\
DecisionTree           & 0.380 & 0.510 & 0.595 & 0.640 & 0.592 & 0.389 & 0.543 $\pm$ 0.092 & 0.396 & 5.00 \\
LogisticRegression ($\ell_2$) & 0.306 & 0.416 & 0.534 & 0.574 & 0.497 & 0.320 & 0.466 $\pm$ 0.095 & 0.291 & 6.33 \\
LogisticRegression     & 0.307 & 0.416 & 0.534 & 0.574 & 0.497 & 0.320 & 0.465 $\pm$ 0.095 & 0.291 & 6.67 \\
\bottomrule
\end{tabular}}
\end{table}

Table~\ref{tab:cvfolds} is what the Phase 2 evaluation asked for. Read down a
column, the windows differ far more than the models: every candidate is at its
worst on the first fold (early 2017, few orders, a 20\% low-review rate) and at
its best on the fourth. Read across a row, the three tree-based rungs are within
0.0004 of one another on the cross-validated mean against a fold-to-fold
standard deviation of 0.09; the stacked ensemble is ahead of the single boosted
model on three of the six windows and behind on three. Cross-validation cannot
order them, which is why the rank rule was fixed before any test score was seen
and why the deployed choice is reported as a tie resolved by that rule rather
than as a win.

\section{Lateness anatomy: where the late deliveries came from}
\label{app:lateness}

Phase 1 read the monthly late-rate series and attributed its spikes to Black
Friday and to the May 2018 truckers' strike. The Phase 1 feedback pointed out
that the figure does not support the second claim --- the largest peaks are
February and March 2018, and May is unremarkable --- and asked for a source and
a more careful link. \code{src/lateness\_anatomy.py} is that analysis, at three
resolutions, and Figure~\ref{fig:lateness} draws it.

\paragraph{Weekly, not monthly.} A ten-day event cannot be resolved by a
monthly series. In Monday-to-Sunday purchase weeks the strike of 21 May to
1 June 2018 \citep{bloomberg2018strike} --- nationwide road blockades, during
which Correios suspended its time-definite Sedex services and added five
business days to standard parcels \citep{tecnoblog2018correios} --- is plainly
visible, but one week \emph{earlier} than a naive reading expects: orders placed
on 14--20 May were promised 19.5 days under normal conditions and were in
transit when the roads closed, and 16.0\% of them arrived late against a 5.0\%
average over the preceding five weeks. Of all orders in transit during the
strike (4{,}228), 16.9\% were late, against 0.4\% for the May orders delivered
before it (Table~\ref{tab:strike}). What the monthly series records instead is
the platform's reaction: order volume in the strike week halved, the promised
window for orders placed that week rose to 25.3 days and in the following week
to 36.8, and the late rate of those orders was 5.5\% and 0.4\%. That is why June
2018, at 1.4\%, is the quietest month in the data.

\paragraph{February--March 2018 is a different mechanism.} Panel (b) of
Figure~\ref{fig:lateness} splits lead time at the carrier handoff. Outside the
peak, sellers hand parcels to the carrier in about 3.2 days and the carrier
takes 8.7 more; in February--March 2018 the seller's share was unchanged at
3.4 days while the carrier's rose to 13.2, in every customer region. At the
same time the promised window was \emph{tightened}, from 26.3 days in January
to 22.0 in March --- within the S\~{a}o Paulo to S\~{a}o Paulo route alone, from
20.1 to 15.2 days --- so the slack that normally absorbs a slow fortnight was
removed exactly when it was needed. The late rate peaked at 29.0\% in the week
of 26 February, the promise was lengthened again from the week of 9 April, and
by that week the late rate was back under 5\%. The data do not identify the
external cause of the carrier slowdown and we do not assert one; what they do
establish is that it was carrier-side, network-wide and compounded by a
promising decision, which is the careful form of the claim Phase 1 made loosely.

\begin{table}[htbp]
\centering
\caption{The cohorts around the strike, from
\code{report/results/lateness\_strike.csv}. Carrier days run from the carrier
handoff to the doorstep; the promise is the window quoted at checkout.}
\label{tab:strike}
\footnotesize
\begin{tabular}{@{}lrrrr@{}}
\toprule
\textbf{Cohort} & \textbf{orders} & \textbf{late} & \textbf{carrier days} & \textbf{promise, days} \\
\midrule
Bought 1--20 May, delivered before the strike & 2{,}749 & 0.4\%  & 4.8  & 19.5 \\
In transit during the strike                  & 4{,}228 & 16.9\% & 13.1 & 25.5 \\
Bought during the strike, 21 May--1 June      & 1{,}502 & 3.5\%  & 9.2  & 29.4 \\
Bought 2--30 June                             & 5{,}912 & 1.4\%  & 6.7  & 27.5 \\
\bottomrule
\end{tabular}
\end{table}

\clearpage
\section{Responses to the Phase 1 and Phase 2 feedback}
\label{app:feedback}

Both rounds of evaluator feedback were treated as specifications.
Tables~\ref{tab:feedback1} and \ref{tab:feedback2} list each point and the
section, table or figure of this report that acts on it.

\begin{center}
\begin{minipage}{\linewidth}
\centering
\captionof{table}{Phase 1 feedback (on the literature review, the exploratory
analysis and the dashboard) and where this report acts on it.}
\label{tab:feedback1}
\small
\begin{tabular}{@{}p{4.4cm}p{10.1cm}@{}}
\toprule
\textbf{Feedback} & \textbf{Response} \\
\midrule
Survey-based perceived quality and platform ratings were treated as one construct &
Sections~\ref{sec:lit} and \ref{sec:problem} separate them, cite the
self-selection literature \citep{hu2009jshaped, hu2017selfselection,
dellarocas2003digitization}, and define the target as the platform rating;
Section~\ref{sec:limitations} carries the consequence \\[2pt]
Model choices cited original algorithm papers rather than e-commerce
applications; no comparison within the tree family; HistGradientBoosting not
justified &
Section~\ref{sec:lit} argues the rungs from the empirical comparisons and from
e-commerce applications, including three studies on this dataset, and states
why HistGradientBoosting; the ladder results of Section~\ref{sec:results} are
the within-family comparison \\[2pt]
Features not supported by literature &
Section~\ref{sec:features} cites the support for each feature group and uses
the correlation matrix as the empirical screen; Appendix~\ref{app:features}
gives a rationale per feature \\[2pt]
Problem 2 should be evaluated under limited capacity, with Precision@$K$ and
Recall@$K$ &
Section~\ref{sec:results-capacity}, Table~\ref{tab:capacity} and
Figure~\ref{fig:capacity}; the operating point in Section~\ref{sec:cost} is a
capacity; the application returns a top-$K$ triage list and replays any test
day \\[2pt]
Figure 3 attributed the late-rate peaks to the May 2018 strike; the strike
needed a source &
Section~\ref{sec:eda-lateness}, Appendix~\ref{app:lateness} and
Figure~\ref{fig:lateness}: the peaks are February--March 2018; the strike is
cited \citep{bloomberg2018strike, tecnoblog2018correios}, shown at weekly
resolution, and its effect traced through the promised window \\[2pt]
Figure 5a omitted S\~{a}o Paulo; ``roughly triples''; Figure 8 not referenced &
Figure~\ref{fig:geography}a redrawn (the notebook had taken the first eight
states alphabetically); growth is about eight-fold
(Section~\ref{sec:temporal}); the correlation matrix is discussed in
Sections~\ref{sec:eda} and \ref{sec:features} \\[2pt]
Candidate 1 named a stakeholder without saying how the team would use the
prediction &
Table~\ref{tab:levers} and Section~\ref{sec:stakeholders}: four levers ordered
by cost, each with its trigger and the evidence behind it \\[2pt]
Dashboard inaccessible &
Redeployed from the team's current repository under a new URL;
Section~\ref{sec:deployment} explains how it stays reachable \\
\bottomrule
\end{tabular}
\end{minipage}
\end{center}

\begin{center}
\begin{minipage}{\linewidth}
\centering
\captionof{table}{Phase 2 feedback (on the modelling report) and where this
report acts on it.}
\label{tab:feedback2}
\small
\begin{tabular}{@{}p{4.4cm}p{10.1cm}@{}}
\toprule
\textbf{Feedback} & \textbf{Response} \\
\midrule
Report CV scores with their fold-to-fold variation for the Problem 2
candidates &
Table~\ref{tab:review} now carries the mean $\pm$ standard deviation over the
five folds for every candidate (and Table~\ref{tab:reg} for the regression);
Table~\ref{tab:cvfolds} and Figure~\ref{fig:cvfolds} show every candidate on
every window; Section~\ref{sec:results-review} interprets the spread; the
model card shows the same table \\[2pt]
State the deployed Problem 2 model clearly and consistently; justify the
choice &
Section~\ref{sec:results-ensemble} and Table~\ref{tab:deployed}: the
StackingClassifier, as pickled, with the justification; the README, the model
card and this text are tied to the artifact by a validation check \\[2pt]
Refer to Figure 3 in the text; the selected hyperparameter values were said to
be in Appendix C but were not &
Every figure in Appendix~\ref{app:figures} is referred to from the body; the
selected values are in Table~\ref{tab:selected}, generated from the results
tables and the deployed pickles \\[2pt]
The expected benefit assumes that contacting a customer prevents the low
review; state it as a limitation and propose a pilot or A/B test &
Section~\ref{sec:limitations} states it; Section~\ref{sec:pilot},
Appendix~\ref{app:pilot} and Figure~\ref{fig:pilot} design and size the
randomised pilot; the model card carries a calculator for it.
Appendix~\ref{app:timing} adds the part of the assumption the data \emph{can}
check --- whether the customer has already written --- and restates the value
on the reachable share \\
\bottomrule
\end{tabular}
\end{minipage}
\end{center}

\section{Training, validation and test side by side}
\label{app:tvt}

The results tables carry cross-validated, validation and test columns. The
in-sample fit on the training period is the missing fourth view, and it is what
separates over-fitting from drift: a model that has memorised its training
period collapses between training and validation, while a model facing a
changed world drops between validation and test, or in step with the base rate.
\code{src/fit\_diagnostics.py} scores the three deployed artifacts on all three
periods without refitting anything.

\begin{table}[htbp]
\centering
\caption{The deployed models on the training, validation and test periods.
Thresholds and the margin-to-risk mapping are the ones fitted on the validation
period and shipped. Lift is PR-AUC divided by that period's own positive rate.}
\label{tab:tvt}
\footnotesize
{\setlength{\tabcolsep}{2.3pt}
\resizebox{\linewidth}{!}{%
\begin{tabular}{@{}llrrrrrrrr@{}}
\toprule
\textbf{Model} & \textbf{Period} & \textbf{$n$} & \textbf{MAE} & \textbf{RMSE} & \textbf{$R^2$} & & & & \\
\midrule
P1A lead time (Voting)      & training   & 70{,}617 & 4.718 & 7.807 & 0.408 & & & & \\
                            & validation & 12{,}957 & 4.784 & 6.592 & 0.222 & & & & \\
                            & test       & 12{,}630 & 3.898 & 5.342 & 0.009 & & & & \\
\midrule
 & & \textbf{$n$} & \textbf{positives} & \textbf{ROC-AUC} & \textbf{PR-AUC} & \textbf{lift} & \textbf{bal.\ acc} & \textbf{F1} & \textbf{Brier} \\
\midrule
P1B late flag (derived)     & training   & 70{,}617 & 8.8\%  & 0.803 & 0.396 & 4.48$\times$ & 0.713 & 0.391 & 0.068 \\
                            & validation & 12{,}957 & 5.0\%  & 0.787 & 0.166 & 3.34$\times$ & 0.648 & 0.227 & 0.044 \\
                            & test       & 12{,}630 & 7.4\%  & 0.782 & 0.214 & 2.89$\times$ & 0.705 & 0.235 & 0.066 \\
\midrule
P2 low review (Stacking)    & training   & 72{,}380 & 15.8\% & 0.836 & 0.617 & 3.90$\times$ & 0.754 & 0.601 & 0.089 \\
                            & validation & 13{,}094 & 11.4\% & 0.762 & 0.409 & 3.60$\times$ & 0.678 & 0.436 & 0.083 \\
                            & test       & 12{,}855 & 11.0\% & 0.749 & 0.414 & 3.75$\times$ & 0.681 & 0.470 & 0.078 \\
\bottomrule
\end{tabular}}}
\end{table}

Three readings. \textbf{The lead-time blend does not memorise its training
period}: its in-sample MAE, 4.72 days, is level with the validation period's
4.78 and \emph{above} the test window's 3.90, and the fall in $R^2$ from 0.41
in-sample to 0.01 on test is the test period's small variance rather than a
large error --- the level shift of Section~\ref{sec:limitations}, not
over-fitting. \textbf{The derived late flag is stable}: ROC-AUC 0.80, 0.79, 0.78
across the three periods, while its lift over the floor falls from 4.5$\times$
to 2.9$\times$ only because the floor itself moves with the late rate.
\textbf{The stacking ensemble shows the expected, modest in-sample optimism of a
tree ensemble}: ROC-AUC 0.84 on the period it was fitted to against 0.76 and
0.75 on the two it never saw, and PR-AUC 0.62 against 0.41 --- but validation
and test agree with each other (lift 3.6$\times$ and 3.75$\times$, Brier
improving from 0.089 to 0.078 as the base rate falls), which is the comparison
that governs deployment.

\section{Feature catalogue}
\label{app:features}

Every feature, with its definition, the reason it is in the model and the
moment it becomes knowable. Literature pointers are to the works discussed in
Section~\ref{sec:lit}; empirical pointers are to Section~\ref{sec:eda}. All 36
are built by \code{src/features.py} from the nine tables alone.

{\footnotesize
\renewcommand{\code}[1]{\texttt{#1}}
\setlength{\tabcolsep}{3pt}
\begin{longtable}{@{}p{3.3cm}p{3.9cm}p{6.6cm}p{1.2cm}@{}}
\caption{The 36 features by group. ``Checkout'' features are admissible for
both problems; ``delivery'' features only for Problem 2.}
\label{tab:features} \\
\toprule
\textbf{Feature} & \textbf{Definition} & \textbf{Why it is in the model} & \textbf{Known at} \\
\midrule
\endfirsthead
\toprule
\textbf{Feature} & \textbf{Definition} & \textbf{Why it is in the model} & \textbf{Known at} \\
\midrule
\endhead
\bottomrule
\endfoot
\multicolumn{4}{@{}l}{\textit{The promise (2)}} \\
\code{estimated\_days} & promised delivery date minus purchase time, days & the platform's own routing knowledge, shown to the customer at checkout; strongest single correlate of lead time after distance ($r = 0.38$) and a decision variable in its own right \citep{cui2024promise, salari2022delivery} & checkout \\
\code{long\_promise} & promised window above the pinned median of 23.2 days & flags the routes the platform itself expects to be slow; pinned so serving and training use the same cut & checkout \\
\multicolumn{4}{@{}l}{\textit{Basket and money (11)}} \\
\code{n\_items}, \code{n\_sellers} & items and distinct sellers in the order & multi-item and multi-seller orders consolidate before shipping; order composition is the standard control in order-level logistics studies \citep{deshpande2023logistics} & checkout \\
\code{is\_multi\_seller} & more than one seller & split shipments wait for the slowest seller & checkout \\
\code{log\_price} & $\log(1 + \text{item total})$ & money is right-skewed; the log stabilises the linear family; price relates to the rating, not to lead time (Figure~\ref{fig:correlations}) & checkout \\
\code{price\_per\_item} & item total divided by items & separates one expensive item from many cheap ones & checkout \\
\code{total\_freight} & freight paid, R\$ & the platform's own pricing of distance and bulk & checkout \\
\code{freight\_ratio} & freight divided by item total & the delivery cost as the customer experiences it, relative to the purchase & checkout \\
\code{log\_weight}, \code{max\_volume\_cm3} & $\log(1 + \text{total weight})$; bulkiest item's volume & handling and carrier constraints; physical attributes are standard delivery-time inputs \citep{salari2022delivery} & checkout \\
\code{total\_photos} & listing photos across the items & a listing-quality proxy for the expectation gap behind a low rating & checkout \\
\code{max\_installments} & largest instalment count & payment behaviour and an affordability signal characteristic of this market & checkout \\
\multicolumn{4}{@{}l}{\textit{Calendar (6)}} \\
\code{purchase\_month}, \code{purchase\_dow}, \code{purchase\_hour} & month, weekday, hour of purchase & seasonality and the weekly and daily rhythm of fulfilment (Figure~\ref{fig:rhythm}); the usual calendar encodings for transactional data \citep{kuhn2019feature, heaton2016feature} & checkout \\
\code{is\_weekend}, \code{is\_business\_hours} & weekend; 09:00--18:00 & an order placed outside working hours waits for the next handling day & checkout \\
\code{is\_nov\_peak} & purchased on or after 20 November & the Black Friday surge stresses fulfilment: the late rate reached 14.3\% in November 2017 (Section~\ref{sec:eda-lateness}) & checkout \\
\multicolumn{4}{@{}l}{\textit{Geography (4)}} \\
\code{log\_distance} & $\log(1 + \text{km})$ between postcode centroids & the strongest checkout-time correlate of lead time ($r = 0.39$); distance is the backbone of delivery-time models \citep{liu2021lastmile} & checkout \\
\code{same\_state} & seller and customer in the same state & the cheapest route class; only 35.7\% of orders & checkout \\
\code{customer\_state} & destination state (27 levels) & regional logistics; rare levels pooled by the encoder & checkout \\
\code{seller\_region} & origin region (5 levels) & origin logistics, collapsed from 27 states so the linear family is not swamped by sparse levels & checkout \\
\multicolumn{4}{@{}l}{\textit{Catalogue and payment (2)}} \\
\code{lead\_category} & category of the most expensive item (74 levels, rare ones pooled) & product type drives both handling and expectations; the category's effect is otherwise invisible to the model & checkout \\
\code{lead\_payment\_type} & payment method of the largest payment & a boleto is paid at a bank and clears days after a card, so handling starts later; the top permutation feature for lead time (Section~\ref{sec:insight}) & checkout \\
\multicolumn{4}{@{}l}{\textit{Track records (8), each an as-of aggregate over orders that had closed before this one was placed}} \\
\code{hist\_seller\_\allowbreak lead\_days}, \code{hist\_seller\_\allowbreak late\_rate}, \code{hist\_seller\_\allowbreak low\_review\_rate} & the seller's mean lead time, late rate and 1--2 star rate over at least five earlier orders & a party's past performance is the standard predictor in supply-chain delay models \citep{baryannis2019risk}; the time-respecting form of target encoding \citep{miccibarreca2001target} & checkout \\
\code{hist\_seller\_n} & number of earlier orders the seller estimate rests on & cold start as information: the model can discount a seller it has barely seen & checkout \\
\code{hist\_route\_\allowbreak lead\_days} & mean lead time on the seller-state to customer-state route & origin--destination history, the route-level analogue of the above \citep{salari2022delivery} & checkout \\
\code{hist\_category\_\allowbreak low\_review\_rate} & the category's past 1--2 star rate & categories differ in how often they disappoint; encoded without seeing the future & checkout \\
\code{hist\_platform\_\allowbreak lead\_days}, \code{hist\_platform\_\allowbreak late\_rate} & platform-wide trailing 28-day mean lead time and late rate & the drift tracker: tells the model how fast the network is running right now \citep{gama2014drift} & checkout \\
\multicolumn{4}{@{}l}{\textit{Delivery outcome (3), Problem 2 only}} \\
\code{delivery\_days}, \code{delay\_vs\_\allowbreak estimate\_days}, \code{is\_late} & realised lead time; days late against the promise; late flag & the dominant driver of a 1--2 star review, 9.2\% to 54.0\% (Section~\ref{sec:reviews}; \citealt{deshpande2023logistics}); known once the parcel has arrived and forbidden before & delivery \\
\end{longtable}}

\section{The outreach pilot}
\label{app:pilot}

\begin{table}[htbp]
\centering
\caption{Pilot size by assumed effect, from \code{report/results/pilot\_design.csv}.
The control-arm rate is the model's precision at that daily capacity on the
test window; $n$ is orders per arm for a two-sided two-proportion test at
$\alpha = 0.05$; duration assumes $K$ contacts a day split evenly across the
arms; net benefit per contact is under the R\$5 : R\$40 assumption.
Figure~\ref{fig:pilot}.}
\label{tab:pilot}
\scriptsize
\setlength{\tabcolsep}{3pt}
\begin{tabular}{@{}rrrrrrrr@{}}
\toprule
\textbf{$K$/day} & \textbf{control} & \textbf{reduction} & \textbf{treated} & \textbf{$n$/arm (80\%)} & \textbf{wks (80\%)} & \textbf{wks (90\%)} & \textbf{net/contact} \\
\midrule
10 & 0.572 & 15\% & 0.486 & 531 & 16 & 21 & $-$R\$1.57 \\
10 & 0.572 & 20\% & 0.457 & 299 & 9  & 12 & $-$R\$0.43 \\
10 & 0.572 & \textbf{21.9\% (break-even)} & 0.447 & 250 & 8 & --- & R\$0.00 \\
10 & 0.572 & 25\% & 0.429 & 191 & 6  & 8  & R\$0.72 \\
10 & 0.572 & 30\% & 0.400 & 133 & 4  & 6  & R\$1.86 \\
10 & 0.572 & 40\% & 0.343 & 74  & 3  & 3  & R\$4.15 \\
\midrule
20 & 0.440 & 20\% & 0.352 & 484 & 7  & 10 & $-$R\$1.48 \\
20 & 0.440 & \textbf{28.4\% (break-even)} & 0.315 & 235 & 4 & --- & R\$0.00 \\
20 & 0.440 & 30\% & 0.308 & 210 & 3  & 5  & R\$0.28 \\
20 & 0.440 & 40\% & 0.264 & 115 & 2  & 3  & R\$2.04 \\
\bottomrule
\end{tabular}
\end{table}

The design randomises inside the top-$K$ list, so both arms share the same
predicted risk and the comparison is a fair one. Two things follow from the
table that a verbal recommendation would miss. The break-even effect grows with
capacity --- 21.9\% at ten contacts a day, 28.4\% at twenty, 42.1\% at forty ---
because precision falls as the list is worked deeper, so a contact lower down
has less to prevent; a pilot should therefore run at the capacity that would
actually be deployed. And the pilot is short at the effect sizes that matter:
six weeks detects a one-in-four reduction at ten contacts a day, while an effect
too small to pay for itself would take sixteen weeks to confirm, which is itself
informative --- if the pilot has not found an effect in two months, the policy
should not be scaled.

\section{When is the review written?}
\label{app:timing}

Problem 2 is framed at delivery, before the review is written. Olist's review
table carries two timestamps --- \code{review\_creation\_date}, the day the
satisfaction survey was sent, and \code{review\_answer\_timestamp}, the moment
the customer answered --- and comparing the second with the delivery timestamp
tests that framing (\code{src/review\_timing.py}). It holds for parcels that
arrive on time and fails for most late ones. The survey goes out at delivery or,
for a parcel still on the road, once the promised date has passed, so the
customer of a late parcel is usually asked to review while still waiting
(Table~\ref{tab:timing}).

\begin{table}[htbp]
\centering
\caption{Review timing on the 95{,}561 delivered and reviewed orders of the
analysis window, from \code{report/results/review\_timing.csv}. A negative median
means the review was answered before the parcel arrived.}
\label{tab:timing}
\small
\begin{tabular}{@{}lrr@{}}
\toprule
 & \textbf{Late parcels} & \textbf{On-time parcels} \\
\midrule
Delivered and reviewed orders & 7{,}659 & 87{,}902 \\
1--2 star rate & 54.0\% & 9.2\% \\
Survey sent on or after the promised date & 99.0\% & 2.1\% \\
1--2 star reviews written before the parcel arrived & \textbf{87.3\%} & 0.5\% \\
Share of all 1--2 star reviews & 33.8\% & 66.2\% \\
Median answer of a 1--2 star review, days after delivery & $-$4.5 & 2.1 \\
\bottomrule
\end{tabular}
\end{table}

\begin{table}[htbp]
\centering
\caption{The review-risk list on the test window, its hits split into customers
who had not yet written (reachable) and those who already had, from
\code{report/results/review\_timing\_capacity.csv}. Precision is the share of
contacts that reach a 1--2 star reviewer; the two columns after it sum to it.}
\label{tab:timing-capacity}
\small
\begin{tabular}{@{}lrrrr@{}}
\toprule
\textbf{Capacity} & \textbf{precision} & \textbf{reachable} & \textbf{already written} & \textbf{late share of contacts} \\
\midrule
top 1\%  & 0.705 & 0.419 & 0.287 & 0.403 \\
top 5\%  & 0.678 & 0.362 & 0.316 & 0.521 \\
top 10\% & 0.475 & 0.302 & 0.173 & 0.376 \\
5 a day  & 0.631 & 0.353 & 0.278 & 0.434 \\
10 a day & \textbf{0.572} & \textbf{0.352} & 0.220 & 0.362 \\
20 a day & 0.440 & 0.288 & 0.152 & 0.298 \\
\midrule
contacting at random & 0.110 & 0.093 & 0.017 & --- \\
\bottomrule
\end{tabular}
\end{table}

Two consequences follow. \textbf{A hit is not always a review that can still be
prevented.} At ten contacts a day, 0.572 of contacts reach a would-be low
reviewer, but 0.220 reach one who has already written --- almost always on a late
parcel --- and 0.352 one who has not, against 0.093 at random
(Table~\ref{tab:timing-capacity}). Counting only the reachable, the expected
benefit under R\$5 : R\$40 falls from R\$1{,}466 to \textbf{R\$775 per thousand
delivered orders} at the same optimum (threshold 0.11), and at ten contacts a
day from R\$853 to R\$433. It stays positive: reachable precision is still almost
three times the 0.125 break-even. \textbf{The action for a late parcel belongs on
the promised date, not at delivery.} Customer care can reach the customer as the
promise breaks, with an apology and a new date, and the survey can be sent after
the parcel arrives. That contact needs no model, since the calendar identifies
the parcels, and it is tested as the second arm of the pilot
(Section~\ref{sec:pilot}). The review-risk list keeps its value for the parcels
that arrive on time, which produce two thirds of all low reviews and are reviewed
a median 2.1 days after delivery. On those parcels alone the deployed model
reaches ROC-AUC 0.711, and its top 5\% are 53.8\% low reviewers, 5.7 times their
base rate, so the list does more than restate lateness
(\code{review\_timing\_subsets.csv}).

The order that opens the presentation shows the pattern in one case. Bought on 16
July 2018 in São Gonçalo (RJ) from a seller in São Luís (MA), 2{,}267\,km away,
and promised in 21 days, it was surveyed two days after the promise passed and
rated one star 5.8 days before it arrived, 10.5 days late. It tops the 16 July
review-risk list at 0.85; eight of that day's top ten left a 1--2 star review,
four of them before their parcel arrived (\code{story\_order.csv}).

\section{Contributions}
\label{app:contributions}

\begin{center}
\small
\begin{tabular}{@{}lp{11.2cm}@{}}
\toprule
\textbf{Member} & \textbf{Contribution} \\
\midrule
Pan Pinyou & Problem scoping and stakeholder framing; feature engineering for the
track-record features; business interpretation, the cost analysis and the pilot
design; report Sections~\ref{sec:problem}, \ref{sec:features} and
\ref{sec:insight}; presentation slides \\[3pt]
Li Yudan & Literature review and its revision; modelling pipeline and the family
ladders; hyperparameter search; imbalance, calibration and fold-to-fold
analysis; report Sections~\ref{sec:lit}, \ref{sec:method} and
\ref{sec:results-review} \\[3pt]
Zhu Gengyu & Data pipeline and the leakage guard; exploratory analysis and the
lateness anatomy; cross-validation protocol and the direct-against-derived
comparison; dashboard, deployment, tests and reproducibility; report
Sections~\ref{sec:data}, \ref{sec:eda}, \ref{sec:results-reg},
\ref{sec:deployment} and \ref{sec:limitations} \\
\bottomrule
\end{tabular}
\end{center}

All three members reviewed every section and verified the reported numbers
against \code{report/stats\_phase3.txt} and \code{report/stats.txt}.
